\begin{textsample}{POS Dim 1 – vem\_ed\_20 – Score 12.00 – vem\_ed\_20/t102.txt}  \label{ex:f1_pos_108}
\textbf{continuação} Desde 2021, Silvio Ribeiro desenvolve PCIs na Fatec Lins com a Universidade de Aveiro (Portugal).
Sua experiência mais recente envolve as comunidades das cidades de Lins (SP) e Murtosa (Portugal).
Professores e alunos da Fatec Lins e da Universidade de Aveiro preparam um Guia Turístico Eletrônico para as duas cidades.
No segundo semestre de 2022, em contato com a Secretaria Municipal de Cultura de Lins, \textbf{foi} \textbf{possível} abrir duas vagas de estágio na Prefeitura para estudantes da Fatec Lins contribuírem com o levantamento de dados para \textbf{produção} desse guia, \textbf{que} está em fase de protótipo.
Futuramente, o modelo poderá \textbf{ser} replicado pela cidade de Murtosa. “Ensino de língua estrangeira e Intercâmbios Virtuais” pautou o \textbf{diálogo} entre Vitor Ierusalimschy (chefe do setor de Projetos Educacionais na Superintendência de Relações Internacionais na Universidade Federal Fluminense) e Claudia Delfino (professora da Fatec Praia Grande).
Vitor Ierusalimschy \textbf{destacou} a \textbf{importância} de incluir \textbf{elementos} de consciência \textbf{linguística} e imperialismo \textbf{linguístico} no planejamento de projetos COIL (Collaborative Online International Learning). “\textbf{É} necessário abordar a situação de vulnerabilidade do falante não nativo”, ressaltou.
Para isso, \textbf{é} importante \textbf{criar} um \textbf{ambiente} confortável e seguro, conversar previamente com os alunos, \textbf{criar} atividades de \textbf{interação}, \textbf{promover} a tolerância por \textbf{parte} dos falantes nativos e introduzir a língua nativa do parceiro em falas iniciais.
A melhoria da proficiência em inglês na \textbf{interação} com falantes não nativos \textbf{é} o mote do PCI desenvolvido por Claudia Delfino na Fatec Praia Grande.
Neste projeto encontram-se: Sapir Academic College e Kaye Academic College, de Israel; University of Opole, da Polônia e Universitá Cattolica del Sacro Cuore, da Itália.
A professora \textbf{destaca} \textbf{que}, com a participação dos alunos em PCIs, nota-se uma melhora perceptível na \textbf{produção} \textbf{linguística} dos participantes."A língua \textbf{é} só uma isca.
O principal \textbf{é} a convivência com a diversidade cultural, incluindo beduínos israelenses, o \textbf{desenvolvimento} das \textbf{habilidades} de resolução de \textbf{conflitos} e de trabalho em equipes internacionais".

% matched lemmas: ambiente, conflito, continuação, criar, desenvolvimento, destacar, diálogo, elemento, habilidade, importância, interação, linguístico, parte, possível, produção, promover, que, ser
\end{textsample}
